\documentclass[10pt,a4paper]{amsart}
\usepackage[margin=19mm]{geometry}
\usepackage{amsmath,amssymb,mathtools}
\usepackage[scr=rsfs,cal=euler]{mathalfa}
\usepackage{xfrac}
\usepackage[unicode,hidelinks,bookmarks=false]{hyperref}
\newcommand{\ie}{i.e.\ }
\newcommand{\cf}{cf.\ }
\newcommand{\resp}{respectively\ }
\newcommand{\Subsection}{\subsection}
\providecommand{\nobreakdash}{\nobreak\hbox{-}}
\setlength{\emergencystretch}{2em}
\multlinegap0pt
\usepackage[shortlabels]{enumitem}
\usepackage{booktabs}
\usepackage{bm}
\usepackage{natbib}
\usepackage{amssymb}
\usepackage{mathtools}
\usepackage{algorithm}\usepackage{algpseudocode}
\usepackage{float}
\newtheorem{lemma}{Lemma}
\newcommand{\HH}{X}
\newcommand{\R}{\mathbb{R}}
\newcommand{\W}{W}
\DeclareMathOperator*{\argmin}{argmin}
\newcommand{\norm}[1]{\lVert{#1}\rVert}
\newcommand{\pair}[2]{{\langle{{#1},{#2}}\rangle}}
\DeclareMathOperator{\prox}{prox}
\newcommand{\reg}{\varphi}
\newcommand{\scalarp}[1]{\langle{#1}\rangle}
\newcommand{\uu}{g}
\newcommand{\x}{x}
\newcommand{\y}{y}
\title{AdaGrad-Diff: A New Version of the Adaptive Gradient Algorithm}
\author{Matia Bojovic, Saverio Salzo, Massimiliano Pontil}
\date{Source: arXiv:2602.13112v1 (CC BY 4.0)}
\begin{document}
\maketitle

\noindent\emph{Source and licence.} AdaGrad-Diff: A New Version of the Adaptive Gradient Algorithm. Version arXiv:2602.13112v1 (CC BY 4.0), distributed by arXiv under the Creative Commons Attribution 4.0 International licence, http://creativecommons.org/licenses/by/4.0/, as recorded in the arXiv licence metadata for this exact version and shown on the abstract page https://arxiv.org/abs/2602.13112v1. Full licence notice: \texttt{licenses/arxiv-2602.13112-CC-BY-4.0.txt}.\par\smallskip

\noindent\textit{Editorial scope.} Counted unit: the article's lemma lem:basicPGstep (source0.tex lines 841--867). The article's complete original proof of it is delivered with it (source0.tex lines 868--941, one \texttt{proof} environment of 74 lines). No mathematics is written, repaired or re-derived: the statement and the proof are the article's own lines, in the article's own notation, and no retained line is substituted or re-flowed.  No further source-local context is needed: every cross-reference of every retained line resolves inside the two delivered spans.  Nothing was omitted from any retained span, so the package carries no omission record.  No retained line contains a citation, so the wrapper carries no bibliography and no bibliography entry.  No retained line contains a figure environment, an image-inclusion command or drawing code, so no image is delivered and the package declares no asset dependency.  Editorial formatting only, in addition: an AMS \texttt{amsart} preamble that restates the article's own declarations and macros used by the retained lines, namely the environments \texttt{lemma} on separate counters, together with the standard packages those lines need.  The retained environments print in this document as Lemma 1 in order of appearance, whereas the article prints the same items with the numbering of its own full document.  The complete native source of the article as distributed in the arXiv e-print for this exact version is bundled unchanged as \texttt{sources/194617/source0.tex}.
\begin{lemma}
\label{lem:basicPGstep}
Let $X$ be a real Hilbert space,
let $\W\colon \HH\to \HH$ be a positive self-adjoint operator bounded from below and set
$\scalarp{\x,\y}_{\W} = \scalarp{\W\x,\y}$, for every $\x,\y\in \HH$.
Let $f,\reg\in \Gamma_0(X)$ and set $F = f+\reg$.
Let $x \in X$, $\uu \in \partial f(x)$, $\eta>0$, and set 
%$\x_+ = \argmin_{\y\in \HH} \big\{ \scalarp{\y-\x, \uu} + \reg(\y)+ (2\eta)^{-1}\norm{\y-\x}_{\W}^2 \big\}$.
$x_+ = \prox^\W_{\eta g}(x - \eta \W^{-1}\uu) = \argmin_{\y\in \HH} \big\{ \scalarp{\y-\x, \uu} + \reg(\y)+ (2\eta)^{-1}\norm{\y-\x}_{\W}^2 \big\}$. 
Then, for every $z\in \HH$
\begin{equation}
\label{eq:basicPGstep}
2\eta (F(x_+)- F(z)) \leq \norm{x-z}_{\W}^2 - \norm{x_+-z}_{\W}^2 - \norm{x_+-x}_{\W}^2 +2\eta \big[f(x_+)-f(x)- \pair{x_+-x}{\uu}\big].
\end{equation}
%In particular the following hold.
%\begin{enumerate}[label={\rm (\roman*)}]
%\item\label{lem:basicPGstep_i} Suppose that $\uu_+\in \partial f(x_+)$. Then
%\begin{equation*}
%(\forall\,z\in X)\quad 2\eta (F(x_+)- F(z)) \leq \norm{x-z}_{\W}^2 - \norm{x_+-z}_{\W}^2 +\eta^2 \norm{\uu - \uu_+}_{\W^{-1}}^2.
%\end{equation*}
%\item\label{lem:basicPGstep_ii} Suppose that $f\colon X\to \R$ is Lipschitz smooth with constant $L$. Then
%\begin{equation*}
%(\forall\,z\in X)\quad 2\eta (F(x_+)- F(z)) \leq \norm{x-z}_{\W}^2 - \norm{x_+-z}_{\W}^2 + \eta^2 \norm{\nabla f(x_+)-\nabla f(x)}^2
%- \frac{\eta}{L} \norm{\nabla f(x)-\nabla f(x_+)}^2.
%\end{equation*}
%\end{enumerate}
\end{lemma}

\begin{proof}
%As a preliminary observation we note that $(\HH, \scalarp{\,\cdot\,,\,\cdot\,}_\W)$ is also a Hilbert space.
%We first prove \eqref{eq:basicPGstep}. 
Let $z\in X$.
By definition of $x_+$ and Fermat's rule, $0\in \uu + \W(\x_+-\x)/\eta +  \partial \reg(\x_+)$
we have $\W(x-x_+)/\eta - \uu\in \partial \reg(x_+)$. Thus, by definition of subdifferential,
\begin{equation*}
\reg(z)-\reg(x_+) \geq \pair{z-x_+}{\W(x-x_+)/\eta - \uu} = \eta^{-1}\pair{z-x_+}{\W(x-x_+)}
-\pair{z-x_+}{\uu}.
\end{equation*}
Multiplying both side by $-2\eta$ we have
\begin{equation*}
2\eta(\reg(x_+)-\reg(z)) \leq  2\pair{x_+-z}{\W(x-x_+)}
+2\eta\pair{z-x_+}{\uu}.
\end{equation*}
Now we note that
\begin{equation*}
\norm{\x- z}_{\W}^2 = \norm{\x - x_+}_{\W}^2 + \norm{\x_+-z}_{\W}^2 +2\pair{\W(\x-x_+)}{\x_+-z}.
\end{equation*}
Hence
\begin{equation}
\label{eq:20250829a}
2\eta(\reg(\x_+)-\reg(z)) \leq  \norm{\x- z}_{\W}^2 - \norm{\x_+-z}_{\W}^2 -\norm{\x - \x_+}_{\W}^2 
+2\eta\pair{z-x_+}{\uu}.
\end{equation}
On the other hand, using the fact that $\uu \in \partial f(\x)$, it holds
\begin{equation*}
\pair{z-\x_+}{\uu} = \pair{z-\x}{\uu} + \pair{\x-\x_+}{\uu}
\leq f(z)-f(\x_+) + f(\x_+)- f(\x) - \pair{\x_+-\x}{\uu}.
\end{equation*}
Thus, 
\begin{equation}
\label{eq:20250829b}
2\eta \pair{z-\x_+}{\uu} \leq 2\eta (f(z)-f(\x_+)) + 2 \eta \big[f(\x_+)- f(\x) - \pair{\x_+-\x}{\uu}\big].
\end{equation}
Combining \eqref{eq:20250829a} and \eqref{eq:20250829b}, inequality \eqref{eq:basicPGstep} follows.
%
%\ref{lem:basicPGstep_i}: Referring to inequality \eqref{eq:basicPGstep}, we note that, being $\uu_+\in \partial f(x_+)$,
%\begin{equation*}
%f(x_+)-f(x)- \pair{x_+-x}{u} \leq \pair{x_+-x}{\uu_+}-\pair{x_+-x}{\uu} = \pair{x_+-x}{\uu_+-\uu}.
%\end{equation*}
%Thus, 
%\begin{equation*}
%(\forall\,z\in X)\quad 2\eta (F(x_+)- F(z)) \leq \norm{x-z}_{\W}^2 - \norm{x_+-z}_{\W}^2 - \norm{x-x_+}_{\W}^2 +2\eta \pair{x-x_+}{\uu - \uu_+}.
%\end{equation*}
%On the other hand using Young-Fenchel inequality we have
%\begin{align*}
%2\eta \pair{x-x_+}{\uu - \uu_+}-\norm{x-x_+}_{\W}^2
%&= 2\eta \pair{x-x_+}{\W^{-1}(\uu - \uu_+)}_{\W}-\norm{x-x_+}_{\W}^2\\[1ex]
%& =2 \bigg(\pair{x-x_+}{\eta\W^{-1}(\uu - \uu_+)}_{\W} -\frac{\norm{x-x_+}_{\W}^2}{2}\bigg)\\
%&\leq 2 \frac{\norm{\eta\W^{-1}(\uu -\uu_+)}_{\W}^2}{2} = \eta^2\norm{\uu- \uu_+}_{\W^{-1}}^2
%\end{align*}
%and the statement follows.
%
%\ref{lem:basicPGstep_ii}: 
%By the Lipschitz smoothness of $f$ we have
%\begin{equation*}
%\frac{1}{2L}\norm{\nabla f(x_+)-f(x)}^2\leq f(x)-f(x_+)- \pair{x-x_+}{\nabla f(x_+)}. 
%\end{equation*}
%Thus, from \eqref{eq:basicPGstep} (with $\uu= \nabla f(\x)$) we have
%\begin{align*}
%2\eta (F(x_+)- F(z)) &\leq \norm{x-z}_{\W}^2 - \norm{x_+-z}_{\W}^2 - \norm{x_+-x}_{\W}^2 +2\eta \big[f(x_+)-f(x)\big] - 2\eta\pair{x_+-x}{\nabla f(x)}\\[1ex]
%&\leq\norm{x-z}_{\W}^2 - \norm{x_+-z}_{\W}^2 - \norm{x_+-x}_{\W}^2\\
%&\qquad +2\eta \Big[\pair{x_+-x}{\nabla f(x_+)}- \frac{1}{2L} \norm{\nabla f(x)-\nabla f(x_+)}^2\Big] - 2\eta\pair{x_+-x}{\nabla f(x)}\\[1ex]
%& =\norm{x-z}_{\W}^2 - \norm{x_+-z}_{\W}^2\\
%&\hspace{15ex} + 2\Big[\pair{x_+-x}{\eta\W^{-1}(\nabla f(x_+)-\nabla f(x))}_{\W} - \frac1 2 \norm{x_+-x}_{\W}^2\Big]
%- \frac{\eta}{L} \norm{\nabla f(x)-\nabla f(x_+)}^2\\[1ex]
%&\leq \norm{x-z}_{\W}^2 - \norm{x_+-z}_{\W}^2 +  \norm{\eta\W^{-1}(\nabla f(x_+)-\nabla f(x))}_{\W}^2
%- \frac{\eta}{L} \norm{\nabla f(x)-\nabla f(x_+)}^2\\[1ex]
%&\leq \norm{x-z}_{\W}^2 - \norm{x_+-z}_{\W}^2 +  \eta^2\norm{\nabla f(x_+)-\nabla f(x)}_{\W^{-1}}^2
%- \frac{\eta}{L} \norm{\nabla f(x)-\nabla f(x_+)}^2.
%\end{align*}
%and the statement follows.
\end{proof}

\end{document}
